% Options for packages loaded elsewhere
\PassOptionsToPackage{unicode}{hyperref}
\PassOptionsToPackage{hyphens}{url}
\documentclass[
]{article}
\usepackage{xcolor}
\usepackage{amsmath,amssymb}
\setcounter{secnumdepth}{5}
\usepackage{iftex}
\ifPDFTeX
  \usepackage[T1]{fontenc}
  \usepackage[utf8]{inputenc}
  \usepackage{textcomp} % provide euro and other symbols
\else % if luatex or xetex
  \usepackage{unicode-math} % this also loads fontspec
  \defaultfontfeatures{Scale=MatchLowercase}
  \defaultfontfeatures[\rmfamily]{Ligatures=TeX,Scale=1}
\fi
% \usepackage{lmodern}
\ifPDFTeX\else
  % xetex/luatex font selection
\fi
% Use upquote if available, for straight quotes in verbatim environments
\IfFileExists{upquote.sty}{\usepackage{upquote}}{}
\IfFileExists{microtype.sty}{% use microtype if available
  \usepackage[]{microtype}
  \UseMicrotypeSet[protrusion]{basicmath} % disable protrusion for tt fonts
}{}
\makeatletter
\@ifundefined{KOMAClassName}{% if non-KOMA class
  \IfFileExists{parskip.sty}{%
    \usepackage{parskip}
  }{% else
    \setlength{\parindent}{0pt}
    \setlength{\parskip}{6pt plus 2pt minus 1pt}}
}{% if KOMA class
  \KOMAoptions{parskip=half}}
\makeatother
\usepackage{color}
\usepackage{fancyvrb}
\newcommand{\VerbBar}{|}
\newcommand{\VERB}{\Verb[commandchars=\\\{\}]}
\DefineVerbatimEnvironment{Highlighting}{Verbatim}{commandchars=\\\{\}}
% Add ',fontsize=\small' for more characters per line
\newenvironment{Shaded}{}{}
\newcommand{\AlertTok}[1]{\textcolor[rgb]{1.00,0.00,0.00}{\textbf{#1}}}
\newcommand{\AnnotationTok}[1]{\textcolor[rgb]{0.38,0.63,0.69}{\textbf{\textit{#1}}}}
\newcommand{\AttributeTok}[1]{\textcolor[rgb]{0.49,0.56,0.16}{#1}}
\newcommand{\BaseNTok}[1]{\textcolor[rgb]{0.25,0.63,0.44}{#1}}
\newcommand{\BuiltInTok}[1]{\textcolor[rgb]{0.00,0.50,0.00}{#1}}
\newcommand{\CharTok}[1]{\textcolor[rgb]{0.25,0.44,0.63}{#1}}
\newcommand{\CommentTok}[1]{\textcolor[rgb]{0.38,0.63,0.69}{\textit{#1}}}
\newcommand{\CommentVarTok}[1]{\textcolor[rgb]{0.38,0.63,0.69}{\textbf{\textit{#1}}}}
\newcommand{\ConstantTok}[1]{\textcolor[rgb]{0.53,0.00,0.00}{#1}}
\newcommand{\ControlFlowTok}[1]{\textcolor[rgb]{0.00,0.44,0.13}{\textbf{#1}}}
\newcommand{\DataTypeTok}[1]{\textcolor[rgb]{0.56,0.13,0.00}{#1}}
\newcommand{\DecValTok}[1]{\textcolor[rgb]{0.25,0.63,0.44}{#1}}
\newcommand{\DocumentationTok}[1]{\textcolor[rgb]{0.73,0.13,0.13}{\textit{#1}}}
\newcommand{\ErrorTok}[1]{\textcolor[rgb]{1.00,0.00,0.00}{\textbf{#1}}}
\newcommand{\ExtensionTok}[1]{#1}
\newcommand{\FloatTok}[1]{\textcolor[rgb]{0.25,0.63,0.44}{#1}}
\newcommand{\FunctionTok}[1]{\textcolor[rgb]{0.02,0.16,0.49}{#1}}
\newcommand{\ImportTok}[1]{\textcolor[rgb]{0.00,0.50,0.00}{\textbf{#1}}}
\newcommand{\InformationTok}[1]{\textcolor[rgb]{0.38,0.63,0.69}{\textbf{\textit{#1}}}}
\newcommand{\KeywordTok}[1]{\textcolor[rgb]{0.00,0.44,0.13}{\textbf{#1}}}
\newcommand{\NormalTok}[1]{#1}
\newcommand{\OperatorTok}[1]{\textcolor[rgb]{0.40,0.40,0.40}{#1}}
\newcommand{\OtherTok}[1]{\textcolor[rgb]{0.00,0.44,0.13}{#1}}
\newcommand{\PreprocessorTok}[1]{\textcolor[rgb]{0.74,0.48,0.00}{#1}}
\newcommand{\RegionMarkerTok}[1]{#1}
\newcommand{\SpecialCharTok}[1]{\textcolor[rgb]{0.25,0.44,0.63}{#1}}
\newcommand{\SpecialStringTok}[1]{\textcolor[rgb]{0.73,0.40,0.53}{#1}}
\newcommand{\StringTok}[1]{\textcolor[rgb]{0.25,0.44,0.63}{#1}}
\newcommand{\VariableTok}[1]{\textcolor[rgb]{0.10,0.09,0.49}{#1}}
\newcommand{\VerbatimStringTok}[1]{\textcolor[rgb]{0.25,0.44,0.63}{#1}}
\newcommand{\WarningTok}[1]{\textcolor[rgb]{0.38,0.63,0.69}{\textbf{\textit{#1}}}}
\usepackage{longtable,booktabs,array}
\newcounter{none} % for unnumbered tables
\usepackage{calc} % for calculating minipage widths
% Correct order of tables after \paragraph or \subparagraph
\usepackage{etoolbox}
\makeatletter
\patchcmd\longtable{\par}{\if@noskipsec\mbox{}\fi\par}{}{}
\makeatother
% Allow footnotes in longtable head/foot
\IfFileExists{footnotehyper.sty}{\usepackage{footnotehyper}}{\usepackage{footnote}}
\makesavenoteenv{longtable}
\setlength{\emergencystretch}{3em} % prevent overfull lines
\providecommand{\tightlist}{%
  \setlength{\itemsep}{0pt}\setlength{\parskip}{0pt}}
\usepackage[]{natbib}
\bibliographystyle{plainnat}
\usepackage{bookmark}
\IfFileExists{xurl.sty}{\usepackage{xurl}}{} % add URL line breaks if available
\urlstyle{same}
\hypersetup{
  colorlinks=true,linkcolor=red,urlcolor=red,citecolor=red,anchorcolor=red,
  pdfcreator={LaTeX via pandoc}}

\author{}
\date{}

\title{COGANT: Codebase-to-GNN Translation\\\normalsize Theory of the Program Graph IR and Practice of the Python Pipeline}
\author{Daniel Ari Friedman\\\footnotesize{Active Inference Institute}\\\footnotesize{\texttt{daniel@activeinference.institute}}\\\footnotesize{\href{https://orcid.org/0000-0001-6232-9096}{ORCID: 0000-0001-6232-9096}}}
\date{\today}

% Core mathematics
\usepackage{amsmath}
\usepackage{amssymb}
\usepackage{amsfonts}
\usepackage{amsthm}

% Document layout
\usepackage{geometry}
\usepackage{float}
\usepackage{graphicx}

% Tables
\usepackage{booktabs}
\usepackage{longtable}
\usepackage{array}

% Algorithm typesetting
\usepackage[ruled,vlined,linesnumbered]{algorithm2e}

% Code listings
\usepackage{listings}

% Typography and formatting
\usepackage{microtype}
\usepackage{xcolor}
\usepackage[binary-units]{siunitx}

% Cross-references and citations
\usepackage{hyperref}
\hypersetup{
    colorlinks=true,
    allcolors=red
}
\usepackage[capitalise,noabbrev]{cleveref}
\usepackage{natbib}

\begin{document}

\maketitle
\thispagestyle{empty}
\tableofcontents
\newpage


\section{Abstract}\label{abstract}

\textbf{COGANT} (Codebase-to-GNN Translation) is a system for turning
software repositories into structured Active Inference artifacts
expressed in the Active Inference Institute's \textbf{Generalized
Notation Notation} (GNN)\footnote{Throughout this manuscript
  \textbf{GNN} refers to the Active Inference Institute's
  \textbf{Generalized Notation Notation}
  (\href{https://github.com/ActiveInferenceInstitute/GeneralizedNotationNotation}{repository}),
  a structured notation for state-space and process models---\emph{not}
  graph neural networks. Downstream consumers, including graph neural
  network training pipelines, can ingest the emitted artifacts, but
  COGANT itself produces Generalized Notation Notation bundles.}. It
occupies the space between classical program analysis---which already
reasons over graphs such as call graphs and code property graphs
\citep{yamaguchi2014modeling}---and data-driven methods that expect
tensors, consistent node and edge vocabularies, and exportable training
bundles
\citep{allamanis2018survey, wu2020comprehensive, scarselli2009graph}.

The design centers on a \textbf{program graph IR}: nodes denote entities
(for example functions and variables), edges denote relations (for
example calls and data flow), and each assertion carries
\textbf{confidence} and \textbf{provenance} so that downstream models
can treat uncertainty explicitly. A \textbf{fixpoint translation engine}
applies declarative rules iteratively until no new semantic mappings
emerge, resolves overlapping mappings by confidence, and reports
coverage gaps explicitly. \textbf{Dynamic enrichment} from coverage
reports and execution traces promotes mappings from a
\texttt{STATIC\_ONLY} tier to \texttt{STATIC\_PLUS\_RUNTIME} when
runtime evidence corroborates static inference. A \textbf{Python
orchestration layer} (session and pipeline APIs, CLI, configuration)
performs discovery, parsing, normalization, graph construction,
translation rules, state-space compilation, validation, and export. A
\textbf{Rust core} (organized as separate crates in the package tree)
holds graph operations, translation, state-space structures, Generalized
Notation Notation formatting, and PyO3 bindings; the authoritative
\textbf{implementation status} distinguishes what is wired in Python
today from staged or planned native acceleration (see
\texttt{../cogant/docs/SPEC.md}).

This manuscript, parallel to the developer documentation under
\href{../cogant/docs/}{\texttt{../cogant/docs/}}, states the theory of
that IR and the practice of running the pipeline, configuring
Generalized Notation Notation exports (canonical GNN Markdown, GNN JSON,
and interop targets such as GraphML, Parquet, and framework-specific
tensor formats), and extending the tool through parsers, rules,
validators, and exporters. It is written to remain valid whether this
tree stays next to the package (\href{../cogant/}{\texttt{../cogant/}})
or is later promoted under
\href{../../../projects/}{\texttt{../../../projects/}} for full template
pipeline rendering.

\newpage

\section{Introduction and Scope}\label{introduction-and-scope}

Machine learning on source code has matured alongside static analysis:
both need structured representations of programs, but learning pipelines
additionally need fixed feature layouts, batchable graphs, and
reproducible exports \citep{allamanis2018survey}. COGANT addresses that
gap by making the path from repository checkout to \textbf{Generalized
Notation Notation (GNN) bundles} explicit, inspectable, and
configurable.

\subsection{Terminology: ``GNN'' in this
manuscript}\label{terminology-gnn-in-this-manuscript}

In this manuscript, \textbf{GNN} refers to the Active Inference
Institute's \textbf{Generalized Notation Notation}
(\href{https://github.com/ActiveInferenceInstitute/GeneralizedNotationNotation}{repository})
\citep{friedman2024gnn}, a structured notation for state-space and
process models---\emph{not} graph neural networks. COGANT translates
source code into this notation, producing program-graph and state-space
artifacts (state space, observation modalities, actions/policies,
likelihood structure, preferences, time settings, parameterization) that
downstream tools can consume. Those downstream tools include graph
neural network training pipelines, which can ingest the exported tensor
views of the program graph, but such neural networks are a
\emph{consumer} of COGANT's output, not its output format. Where the
manuscript discusses graph neural networks as related work, it uses the
phrase ``graph neural networks'' in full; the acronym GNN always refers
to Generalized Notation Notation.

\subsection{Problem}\label{problem}

A research or engineering group typically has:

\begin{itemize}
\tightlist
\item
  \textbf{Source code} in one or more languages.
\item
  \textbf{Analysis goals} that combine classical queries (who calls
  whom?) with learned models (bug detection, retrieval, summarization).
\item
  \textbf{Tooling fragmentation}: compilers and linters produce one
  shape of facts; PyTorch or JAX expect another.
\end{itemize}

COGANT unifies these behind a single pipeline whose intermediate
artifacts are documented as a progression of IRs (repo IR, program
graph, semantic mapping, state space, process model, validation), as
described in \texttt{../cogant/docs/ARCHITECTURE.md} and
\texttt{../cogant/docs/SPEC.md}.

\subsection{What COGANT is not}\label{what-cogant-is-not}

COGANT is not a replacement for a full compiler IR, a theorem prover, or
a security scanner by itself. It \textbf{extracts and serializes}
program graphs and behavioral sketches for downstream use. Language
coverage and rule depth follow the implementation-status table in the
SPEC: Python is the primary front end at v0.1.x; additional languages
and Rust acceleration are staged as documented there.

\subsection{Manuscript versus package
docs}\label{manuscript-versus-package-docs}

The canonical API and CLI details live in:

\begin{itemize}
\tightlist
\item
  \texttt{../cogant/docs/API\_GUIDE.md} --- Session,
  \texttt{PipelineRunner}, \texttt{Bundle}, \texttt{ReviewAPI}
\item
  \texttt{../cogant/docs/CLI\_GUIDE.md} --- commands and flags
\item
  \texttt{../cogant/docs/GNN\_EXPORT.md} --- Generalized Notation
  Notation schema, section ordering, and companion interop field
  contracts (including optional PyG/DGL tensor views)
\item
  \texttt{../cogant/docs/PLUGIN\_API.md} --- parsers, rules, validators,
  exporters
\item
  \texttt{../cogant/docs/ARCHITECTURE.md} --- layers, crates, data flow
\end{itemize}

This manuscript summarizes theory and practice for readers who want a
single linear narrative before diving into those files.

\subsection{Dual Python/Rust
architecture}\label{dual-pythonrust-architecture}

COGANT employs a dual-language architecture. The \textbf{Python
orchestration layer} handles session management, pipeline coordination,
configuration, file discovery, AST parsing, and the plugin interface --
tasks where developer ergonomics and extensibility matter more than raw
throughput. The \textbf{Rust core}, organized as a workspace of seven
crates (\texttt{cogant-core}, \texttt{cogant-graph},
\texttt{cogant-translate}, \texttt{cogant-statespace},
\texttt{cogant-gnn}, \texttt{cogant-store}, \texttt{cogant-ffi}),
implements typed graph operations, translation engine internals,
state-space structures, and Generalized Notation Notation (GNN)
formatting -- tasks where memory layout and cache locality dominate
wall-clock time. Communication between layers flows through PyO3
bindings (\texttt{cogant-ffi}), which release the Python GIL during Rust
execution to permit concurrent stage processing. Where Rust bindings are
not yet wired for a code path (see the implementation-status table in
\texttt{../cogant/docs/SPEC.md}), Python fallback implementations ensure
the pipeline remains functional at the cost of higher latency on large
repositories.

\subsection{Roadmap of the following
sections}\label{roadmap-of-the-following-sections}

Section 2 formalizes the program graph and IR progression. Section 3
describes the realized pipeline behavior and artifacts. Sections 4--7
cover conclusions, how to run experiments, reproducibility, and related
work.

\newpage

\section{Methodology: Program Graphs and Intermediate
Representations}\label{methodology-program-graphs-and-intermediate-representations}

\subsection{Program graph}\label{program-graph}

Let \(G = (V, E)\) denote a directed graph whose vertices \(V\)
represent program entities and whose edges \(E\) represent
relationships. COGANT assigns each node a \textbf{kind} (for example
function, variable, type) and a \textbf{semantic role} (for example
definition versus use). Each node and edge may carry:

\begin{itemize}
\tightlist
\item
  A \textbf{stable identifier} for persistence across runs when hashing
  allows.
\item
  \textbf{Type strings} and attributes recovered from the front end.
\item
  A \textbf{confidence} score in \([0,1]\) and \textbf{provenance}
  metadata (source code, type system, control flow, heuristic, external
  tool).
\end{itemize}

This follows the same philosophical line as code property graphs that
fuse AST, control flow, and data flow into one analyzable structure
\citep{yamaguchi2014modeling}, but orients the representation toward
\textbf{Generalized Notation Notation (GNN) export} and, optionally,
tensorized views of the same graph: node kinds and roles map to discrete
indices in both forms, and optional embeddings can augment features as
described in \texttt{../cogant/docs/GNN\_EXPORT.md}.

\subsubsection{Structural equivalence}\label{structural-equivalence}

Two program graphs \(G_1 = (V_1, E_1)\) and \(G_2 = (V_2, E_2)\) are
usefully compared via typed graph isomorphism: a bijection
\(\phi : V_1 \to V_2\) preserving edge structure and relevant labels
supports deduplication and cross-repository linking when interfaces
align.

\begin{equation}
\label{eq:typed-iso}
(u, v) \in E_1 \iff (\phi(u), \phi(v)) \in E_2
\end{equation}

Equation \ref{eq:typed-iso} formalizes the structural invariant we check
during deduplication and cross-repository linking: a candidate bijection
\(\phi) is accepted only when every edge in \(G_1\) has a corresponding
edge between the images of its endpoints in \(G_2\) (and vice versa), so
that label-preserving matches strictly preserve the adjacency structure
of both program graphs.

\subsection{Progressive IRs}\label{progressive-irs}

Processing proceeds through a sequence of representations (see
\texttt{../cogant/docs/SPEC.md}):

\begin{enumerate}
\def\labelenumi{\arabic{enumi}.}
\tightlist
\item
  \textbf{Repo IR} --- entities and relationships extracted from
  parsers.
\item
  \textbf{Program graph IR} --- consolidated graph with deduplication
  and metadata.
\item
  \textbf{Semantic mapping IR} --- output of the translation rule
  engine.
\item
  \textbf{State space IR} --- variables, actions, transitions,
  observations.
\item
  \textbf{Process model IR} --- higher-level control patterns where
  implemented.
\item
  \textbf{Validation IR} --- coverage, confidence analysis, schema
  checks.
\end{enumerate}

Not every stage is equally complete for every repository; the SPEC marks
partial areas (translation rules, state space, Rust acceleration)
explicitly.

\subsection{Translation rules}\label{translation-rules}

The \textbf{translation} stage applies declarative rules that refine
roles, attach labels, and adjust confidence. Concurrency targets and
layering are described in \texttt{../cogant/docs/ARCHITECTURE.md}. The
rule engine composes passes over the graph in a \textbf{fixpoint loop}:
rules are re-applied until no new semantic mappings emerge or a
configurable iteration cap is reached. Overlapping rules may require
priority ordering as documented there; when multiple rules claim the
same graph fragment, the engine retains the mapping with the highest
confidence score, following the principle that edit representations
should be composable \citep{yin2019learning}.

\subsubsection{Fixpoint iteration, conflict resolution, and
coverage}\label{fixpoint-iteration-conflict-resolution-and-coverage}

The shipped \texttt{TranslationEngine} in
\texttt{../cogant/py/cogant/translate/engine.py} realizes the fixpoint
loop concretely. Each iteration walks every registered rule once: the
engine calls \texttt{rule.matches(graph,\ query)} to collect candidate
fragments, then calls \texttt{rule.apply(graph,\ match)} on each,
accumulating any resulting \texttt{SemanticMapping} objects keyed by
their stable IDs. A per-pass counter tracks mappings that were genuinely
new (not already present in the running set), and the loop terminates as
soon as a pass completes with zero additions. The engine logs each
iteration boundary through an internal match log, so the number of
passes required to reach a fixed point is directly observable in the
post-run diagnostics. The default iteration cap is
\texttt{max\_iterations\ =\ 10}; in testing, most repositories converge
well before that bound, and the cap exists primarily as a safety valve
against pathological rule sets that could otherwise oscillate
indefinitely.

After fixpoint termination, the engine invokes
\texttt{\_resolve\_conflicts()} to reconcile mappings whose
\texttt{graph\_fragment\_node\_ids} sets overlap. For each overlapping
pair the engine retains the mapping with the higher
\texttt{confidence\_score} and discards the other, logging a
\texttt{conflict\_resolved} event that records the losing ID, the
winning ID, and the specific overlap set. Rule priority is therefore
expressed indirectly through the confidence model rather than through a
separate ordering table: a higher-priority rule expresses its precedence
by yielding mappings with higher scores, which then win in the conflict
resolution pass. A companion entry point,
\texttt{translate\_with\_confidence()}, runs the standard fixpoint loop,
rescores every surviving mapping through the \texttt{ConfidenceModel},
and then re-resolves conflicts so that any ordering shifts induced by
rescoring are honored.

Translation coverage -- the fraction of graph nodes that received at
least one semantic mapping -- is reported by
\texttt{get\_coverage\_report(graph)}. It returns the total node count,
the number of covered nodes, the number of uncovered nodes, a
\texttt{coverage\_percent} value rounded to two decimal places, and the
sorted list of uncovered node IDs. The uncovered list is intentionally
emitted verbatim so that downstream tooling can target unmapped regions
for manual review or rule extension, rather than burying the gap behind
a single aggregate number \citep{allamanis2018survey}.

\begin{algorithm}
\caption{Fixpoint translation engine}
\label{alg:fixpoint}
\KwIn{Program graph $G = (V, E)$; rule set $R$; maximum iterations $K$ (default 10)}
\KwOut{Set of semantic mappings $\mathcal{M}$ keyed by stable ID}
$\mathcal{M} \leftarrow \emptyset$\;
\For{$k \leftarrow 1$ \KwTo $K$}{
    $n_{\text{new}} \leftarrow 0$\;
    \ForEach{rule $r \in R$ sorted by $\text{priority}(r)$ descending}{
        \ForEach{match $m \in r.\text{matches}(G)$}{
            $\mu \leftarrow r.\text{apply}(G, m)$\;
            \If{$\mu \neq \bot$ \textbf{and} $\mu.\text{id} \notin \mathcal{M}$}{
                $\mathcal{M} \leftarrow \mathcal{M} \cup \{\mu\}$\;
                $n_{\text{new}} \leftarrow n_{\text{new}} + 1$\;
            }
        }
    }
    \If{$n_{\text{new}} = 0$}{
        \textbf{break} \tcp*{fixed point reached}
    }
}
$\mathcal{M} \leftarrow \textsc{ResolveConflicts}(\mathcal{M})$\;
\Return $\mathcal{M}$\;
\end{algorithm}

Algorithm \ref{alg:fixpoint} summarizes the engine. Termination is
guaranteed because each iteration either produces at least one new
mapping (whose stable ID is then fixed in \(\mathcal{M}\)) or terminates
by the break condition; the outer \(K\) bound serves as a safety valve
for pathological rule sets.

\begin{algorithm}
\caption{Priority-ordered conflict resolution}
\label{alg:conflict}
\KwIn{Mapping set $\mathcal{M}$; priority function $p(\cdot)$}
\KwOut{Reduced mapping set with no overlapping fragments}
Build inverted index $\mathcal{I}: V \to 2^{\mathcal{M}}$ from node IDs to mappings touching them\;
$\mathcal{C} \leftarrow \emptyset$ \tcp*{conflict pairs}
\ForEach{node $v$ with $|\mathcal{I}(v)| \geq 2$}{
    \ForEach{$(\mu_a, \mu_b) \in \binom{\mathcal{I}(v)}{2}$}{
        $\mathcal{C} \leftarrow \mathcal{C} \cup \{(\mu_a, \mu_b)\}$\;
    }
}
$\mathcal{R} \leftarrow \emptyset$ \tcp*{removal set}
\ForEach{$(\mu_a, \mu_b) \in \mathcal{C}$}{
    \If{$\mu_a \in \mathcal{R}$ \textbf{or} $\mu_b \in \mathcal{R}$}{\textbf{continue}}
    $k_a \leftarrow (p(\mu_a), c(\mu_a))$; $k_b \leftarrow (p(\mu_b), c(\mu_b))$\;
    \eIf{$k_a \geq k_b$}{$\mathcal{R} \leftarrow \mathcal{R} \cup \{\mu_b\}$}{$\mathcal{R} \leftarrow \mathcal{R} \cup \{\mu_a\}$}
}
\Return $\mathcal{M} \setminus \mathcal{R}$\;
\end{algorithm}

Algorithm \ref{alg:conflict} detects conflicts via an inverted index in
\(O(\sum_v |\mathcal{I}(v)|^2)\) worst-case time, which is substantially
faster than the naive all-pairs scan for graphs where most nodes carry
at most one mapping.

\subsection{Confidence Scoring and Evidence
Tiers}\label{confidence-scoring-and-evidence-tiers}

The shipped \texttt{ConfidenceModel} in
\texttt{../cogant/py/cogant/translate/confidence.py} computes a scalar
score in \([0,1]\) from:

\begin{itemize}
\tightlist
\item
  Average confidence over provenance records.
\item
  An evidence-diversity term (scaled, capped).
\item
  A parser certainty factor applied multiplicatively.
\item
  Conflict penalties subtracted after scaling.
\end{itemize}

\begin{equation}
\label{eq:confidence-core}
c = \max\left(0,\ \min\left(1,\ (\bar{e} + \delta_d)\cdot \kappa - \pi\right)\right)
\end{equation}

Here \(\bar{e}\) is the mean evidence confidence, \(\delta_d\) is the
diversity bonus (bounded), \(\kappa) is parser certainty, and \(\pi)
aggregates conflict penalties. \textbf{Tiers} (for example static-only
versus static-plus-runtime) are assigned from score thresholds and
evidence source tags (\texttt{determine\_confidence\_tier}); see the
same module for named thresholds and enum values. The manuscript does
not duplicate those literals so they cannot drift from code.

\subsection{State space and behavior}\label{state-space-and-behavior}

Where traces or coverage are available, \textbf{dynamic} extraction
feeds the state-space compiler. The goal is a compact behavioral model:
states, actions, transitions, and observations that sit alongside the
static graph for tasks that require execution-sensitive features.

The shipped \texttt{StateSpaceCompiler} in
\texttt{../cogant/py/cogant/statespace/compiler.py} constructs this
model in several coordinated passes driven by the semantic mappings and
the underlying program graph.

\textbf{State variables} are identified by a
\texttt{StateVariableExtractor} that traverses graph nodes carrying
READS and WRITES edges: any node that participates in a write is a
candidate hidden-state variable, and its type, cardinality, and initial
confidence are derived from the node's recovered type string and the
provenance of the rules that produced the associated mapping.

\textbf{Actions} are extracted from semantic mappings of kind
\texttt{ACTION}, with parameters read from node metadata (typically the
parser-recovered function signature), effects traced through outgoing
WRITES edges on the controller node, and preconditions derived from
parameter lists, docstring directives, and explicit metadata entries.
For method-kind actions, the compiler also walks CONTAINS edges back to
the enclosing class so that instance-level state mutations are
attributed to the correct controller, mirroring how code property graphs
fuse structural and data-flow views \citep{yamaguchi2014modeling}.

\textbf{Transitions} are inferred by a cross-reference pass
(\texttt{\_cross\_reference\_actions\_and\_variables}) that, for each
action, partitions its adjacent variables into reads and writes based on
edge kind (WRITES and MUTATES yielding writes; READS and OBSERVES
yielding reads). The resulting \texttt{Transition} object records a
\texttt{source\_state} in which every touched variable is marked
\texttt{"pre"} and a \texttt{target\_state} in which written variables
advance to \texttt{"post"} while read-only variables remain
\texttt{"pre"}; this simple pre/post convention keeps the model faithful
to the static evidence without overcommitting to symbolic value domains
that cannot be recovered from AST analysis alone. Trigger attribution
follows incoming TRIGGERS and CALLS edges so that orchestration flow is
preserved in the behavioral model.

\textbf{Observation modalities} are built from semantic mappings of kind
\texttt{OBSERVATION}: each associated node becomes an
\texttt{ObservationModality} whose modality type (\texttt{log},
\texttt{metric}, \texttt{event}, \texttt{sensor}, or a generic fallback)
is inferred from the mapping's description and the node's name, giving
the GNN export a typed observation channel aligned with the OBSERVES
edges in the program graph.

The \textbf{temporal regime} of the model is determined by a companion
\texttt{TemporalAnalyzer}. It classifies nodes as asynchronous when
their metadata carries \texttt{is\_async}/\texttt{async} flags or their
names match patterns such as \texttt{async}, \texttt{callback},
\texttt{promise}, or \texttt{future}, and as event-related when their
kind is \texttt{EVENT} or their names match \texttt{event},
\texttt{handler}, \texttt{listener}, or \texttt{trigger}. Temporal
orderings are then extracted from CALLS and TRIGGERS edges, with each
edge classified as \texttt{parallel} (when either endpoint is async) or
\texttt{sequential} otherwise, and event patterns are assembled from
triggers-in / triggers-out pairs around each event node. A final
decision rule selects among \texttt{SYNCHRONOUS}, \texttt{ASYNCHRONOUS},
\texttt{EVENT\_DRIVEN}, and \texttt{HYBRID}: the presence of event
triggers and patterns combined with async handlers yields
\texttt{HYBRID}; event triggers alone yield \texttt{EVENT\_DRIVEN}; an
async fraction above 30 percent or the presence of async handlers yields
\texttt{ASYNCHRONOUS}; and all remaining cases default to
\texttt{SYNCHRONOUS}. This regime is attached to the
\texttt{StateSpaceModel} metadata so downstream consumers know which
execution model the transition graph assumes.

When coverage and traces are available, \texttt{enrich\_graph()} in
\texttt{../cogant/py/cogant/dynamic/enrichment.py} feeds additional
evidence into this pipeline. Coverage enrichment matches
\texttt{.coverage} SQLite databases or Cobertura XML reports against
nodes whose \texttt{path} and \texttt{source\_range} overlap covered
lines, attaching \texttt{coverage\_hits} and, where branch data is
available, \texttt{branch\_coverage} metadata. Trace enrichment parses
Chrome DevTools traces, writes \texttt{call\_count},
\texttt{avg\_duration\_ms}, and \texttt{is\_hot\_path} onto matching
callable nodes, and adds or reweights dynamic CALLS edges tagged with
\texttt{evidence\_sources={[}"dynamic\_trace"{]}}. Both steps also
append \texttt{dynamic\_coverage} and \texttt{dynamic\_trace} markers to
the program graph's evidence sources. The confidence model consumes
these markers directly: any mapping whose evidence set now contains both
static and dynamic entries becomes eligible for promotion from the
\texttt{STATIC\_ONLY} tier to \texttt{STATIC\_PLUS\_RUNTIME}, and
hot-path and branch-coverage metadata raise the underlying score through
the diversity bonus \(\delta_d\) in Equation \ref{eq:confidence-core}.
This is the mechanism by which executing the target program, even
partially, converts static heuristics into corroborated behavioral facts
without rerunning the upstream rule engine \citep{allamanis2018survey}.

\subsection{GNN export (Generalized Notation
Notation)}\label{gnn-export-generalized-notation-notation}

Exports package the program graph and compiled state-space model into
the Active Inference Institute's \textbf{Generalized Notation Notation}
plus a small set of interop targets:

\begin{itemize}
\tightlist
\item
  \textbf{GNN Markdown (\texttt{model.gnn.md})} --- canonical,
  human-readable Generalized Notation Notation organized into the
  section order defined in \texttt{../cogant/docs/GNN\_EXPORT.md} (Model
  Metadata, Repository Metadata, Source Coverage, State Space,
  Observation Modalities, Actions/Policies, Connections, Factors,
  Transition Structure, Likelihood Structure, Preferences/Constraints,
  Time Settings, Parameterization, Ontology Mapping, Provenance,
  Confidence Scores, Rendering Hints, Validation Notes).
\item
  \textbf{GNN JSON (\texttt{model.gnn.json})} --- machine-readable
  equivalent of the markdown sections, plus companion JSON files per
  section (\texttt{state\_space.json}, \texttt{observations.json},
  \texttt{actions.json}, \texttt{transitions.json}, and so on).
\item
  \textbf{Interop exports} --- GraphML and Parquet for analysis in
  Gephi/yEd and DuckDB, and optional tensor views (PyTorch Geometric
  \texttt{Data} objects, DGL graphs, HDF5 tables) for downstream graph
  neural network training pipelines that consume the program graph as a
  relational tensor.
\end{itemize}

Node and edge feature breakdowns, section contracts, and the optional
framework targets are specified in
\texttt{../cogant/docs/GNN\_EXPORT.md}; they determine the structure of
the emitted notation as well as the effective input dimensionality of
any downstream model.

\subsection{Error handling philosophy}\label{error-handling-philosophy}

The implementation distinguishes fatal configuration or parse failures
from per-file errors that allow partial results. Validation aggregates
warnings and inconsistencies into a report that travels with the bundle.
This mirrors the layered error categories documented in the architecture
(fatal, error, warning, info).

\newpage

\section{API and Workflows}\label{api-and-workflows}

This section describes the programmatic surface that the shipped Python
package exposes in practice, aligned with
\texttt{../cogant/docs/API\_GUIDE.md} and
\texttt{../cogant/VERIFICATION\_REPORT.md}. It is \textbf{not} an
empirical benchmark section; COGANT does not ship comparative timing
claims in the manuscript layer. Instead, it records the \textbf{surface
area} users can rely on: two complementary entry points (a Session for
stepwise work and a Pipeline for batch runs), the Bundle accessors that
expose their artifacts, a command-line interface, and a Review API for
human-in-the-loop curation.

\subsection{Session-oriented workflow}\label{session-oriented-workflow}

The Session API is intended for interactive exploration: it is the right
choice when a user wants to inspect intermediate artifacts between
stages, iterate on a single repository in a notebook, or debug a
specific extraction step without committing to a full scripted run. Each
call returns control to the caller so that the graph, mappings, or
state-space model can be examined before the next stage is invoked.

\texttt{Session.from\_target} accepts a local path or URL, then supports
a stepwise workflow:

\begin{itemize}
\tightlist
\item
  \texttt{extract\_static} --- AST-oriented extraction for supported
  languages.
\item
  \texttt{extract\_dynamic} --- traces and coverage when inputs exist.
\item
  \texttt{build\_graph} --- program graph construction.
\item
  \texttt{translate\_to\_gnn} --- Generalized Notation Notation (GNN)
  representation.
\item
  \texttt{compile\_state\_space} --- behavioral model when the pipeline
  has sufficient data.
\item
  \texttt{export\_all} --- writes JSON artifacts under a chosen output
  directory.
\end{itemize}

This path suits interactive notebooks and incremental debugging.

\subsection{Pipeline-oriented
workflow}\label{pipeline-oriented-workflow}

The Pipeline API is intended for scripted, reproducible batch runs: it
is the right choice when a user wants to process many repositories with
a fixed configuration, wire COGANT into CI, or guarantee that every run
executes the same ordered stages with the same plugin settings. Because
the whole run is described by a single \texttt{PipelineConfig}, it can
be checked into version control and replayed without manual
intervention.

\texttt{PipelineRunner} with \texttt{PipelineConfig} runs an ordered
list of stages (ingest, static, normalize, graph, translate, state
space, process, export, validate). Configuration can skip stages (for
example dynamic analysis), attach \textbf{plugin} settings per language,
set \texttt{output\_dir}, verbosity, and dry-run mode. Results aggregate
into a \textbf{Bundle} with \texttt{stage\_results}, error lists, and
accessors described below.

\subsection{Bundle accessors}\label{bundle-accessors}

The bundle API exposes summaries and artifacts, including:

\begin{itemize}
\tightlist
\item
  \texttt{repo\_summary}, \texttt{program\_graph},
  \texttt{state\_space\_model}, \texttt{process\_model}
\item
  \texttt{gnn\_markdown} --- human-readable graph summary
\item
  \texttt{validation\_report}
\item
  \texttt{render\_site} --- static HTML site with graph and model views
\item
  \texttt{to\_json} / \texttt{save\_json}
\end{itemize}

The verification report lists these as delivered capabilities for the
v0.1.0 line.

\subsection{Command-line interface}\label{command-line-interface}

The CLI entry point (\texttt{cogant.cli.main}) implements commands such
as \texttt{init}, \texttt{scan}, \texttt{extract-static},
\texttt{extract-dynamic}, graph and export verbs, and
validation-oriented subcommands. Exact flags live in
\texttt{../cogant/docs/CLI\_GUIDE.md}; the manuscript does not duplicate
them to avoid drift.

\subsection{Review API}\label{review-api}

\texttt{ReviewAPI} supports interactive curation: load a bundle, present
mappings, accept, reject, or edit, then save a curated bundle. This
closes the loop when human review is part of the ML dataset
construction.

\newpage

\section{Examples and Failure Modes}\label{examples-and-failure-modes}

This section walks through a concrete example run of the pipeline on a
small Flask application, shows representative fragments of the artifacts
it produces (semantic mappings and state-space transitions), and then
documents the degradation behavior that users should expect when inputs
are missing or partial. Together these illustrate both what a successful
run looks like and how the system communicates partial success.

\subsection{Example outputs}\label{example-outputs}

The package tree includes examples under
\href{../cogant/examples/}{\texttt{../cogant/examples/}} and sample
exports under
\href{../cogant/output/examples/control_positive/}{\texttt{../cogant/output/examples/control\_positive/}}
(for example \textbf{calculator}, \textbf{event\_pipeline}, and
\textbf{flask\_mini} with \textbf{model.gnn.json} /
\textbf{model.gnn.md} under \texttt{data/} and \texttt{reports/},
diagrams under \texttt{diagrams/}, and raster figures under
\texttt{figures/}). Use \texttt{cogant\ translate\ -\/-layout-output} or
\texttt{PipelineConfig(layout\_output=True)} to place pipeline JSON
under \texttt{data/} automatically; run
\texttt{python\ -m\ cogant.tools.render\_output\_figures} on the output
root for PNGs. Regenerate these trees by running the packaged pipeline
against the example sources when validating releases.

\subsubsection{Concrete walkthrough: Flask REST
API}\label{concrete-walkthrough-flask-rest-api}

To make the pipeline's behavior tangible, consider a small Flask
application with three HTTP endpoints, two utility modules, and one
data-access layer -- 800 lines of Python across 6 files, of which 782
lines are non-blank, non-comment code analyzed by the pipeline. Running
\texttt{cogant\ scan\ -\/-target\ ./flask\_app\ \&\&\ cogant\ extract-static\ \&\&\ cogant\ build-graph\ \&\&\ cogant\ translate\ \&\&\ cogant\ export-all\ -o\ output/}
on this repository produces a program graph with the following
characteristics:

{\def\LTcaptype{none} % do not increment counter
\begin{longtable}[]{@{}ll@{}}
\toprule\noalign{}
Metric & Value \\
\midrule\noalign{}
\endhead
\bottomrule\noalign{}
\endlastfoot
Source files discovered & 6 \\
Lines analyzed & 782 / 800 (97.8\%) \\
\textbf{Nodes} & \textbf{147} \\
\textbf{Edges} & \textbf{389} \\
Mean node confidence & 0.91 \\
Validation status & PASS (0 errors, 3 warnings) \\
\end{longtable}
}

The node and edge populations distribute across kinds as follows:

\textbf{Table 1. Node kind distribution (Flask API example).}

{\def\LTcaptype{none} % do not increment counter
\begin{longtable}[]{@{}lll@{}}
\toprule\noalign{}
Node kind & Count & Percentage \\
\midrule\noalign{}
\endhead
\bottomrule\noalign{}
\endlastfoot
FUNCTION & 38 & 25.9\% \\
VARIABLE & 52 & 35.4\% \\
TYPE & 11 & 7.5\% \\
MODULE & 6 & 4.1\% \\
CONTROLFLOW\_NODE & 18 & 12.2\% \\
DATA\_STRUCTURE & 7 & 4.8\% \\
ERRORHANDLER & 5 & 3.4\% \\
CONSTANT & 6 & 4.1\% \\
EXTERNAL & 4 & 2.7\% \\
\end{longtable}
}

\textbf{Table 2. Edge kind distribution (Flask API example).}

{\def\LTcaptype{none} % do not increment counter
\begin{longtable}[]{@{}lll@{}}
\toprule\noalign{}
Edge kind & Count & Percentage \\
\midrule\noalign{}
\endhead
\bottomrule\noalign{}
\endlastfoot
CALLS & 87 & 22.4\% \\
USES & 104 & 26.7\% \\
DEFINES & 62 & 15.9\% \\
HAS\_TYPE & 48 & 12.3\% \\
DATA\_FLOW & 34 & 8.7\% \\
MEMBER\_OF & 29 & 7.5\% \\
INHERITS & 3 & 0.8\% \\
Other & 22 & 5.7\% \\
\end{longtable}
}

\subsubsection{Example semantic mapping
output}\label{example-semantic-mapping-output}

During the translation stage, each program-graph node is matched against
the active rule set and assigned a semantic role with an associated
confidence score. The following excerpt shows three representative
mappings from the Flask example:

\begin{verbatim}
Node: get_users  (kind=FUNCTION, file=routes/users.py:14)
  Matched rule: rule_fn_def_001
  Target role:  FUNCTION_DEF
  Confidence:   0.98  (base=1.0, -0.02 missing docstring penalty)
  Provenance:   SourceCode

Node: db.session.query  (kind=FUNCTION, file=routes/users.py:22)
  Matched rule: rule_method_call_001
  Target role:  METHOD_CALL
  Confidence:   0.82  (base=0.90, -0.08 receiver type inferred heuristically)
  Provenance:   Heuristic

Node: User  (kind=TYPE, file=models/user.py:5)
  Matched rule: rule_type_def_001
  Target role:  TYPE_DEF
  Confidence:   1.00
  Provenance:   SourceCode
\end{verbatim}

The \texttt{db.session.query} call illustrates how the confidence model
(Equation \ref{eq:confidence-core} in Section 2) penalizes heuristic
provenance: the receiver type of \texttt{session} is resolved by import
tracing rather than explicit annotation, reducing \(\kappa) below 1.0
and yielding a final confidence of 0.82 in the MEDIUM tier.

\subsubsection{Example state-space
excerpt}\label{example-state-space-excerpt}

When dynamic traces are available, the state-space compiler produces a
behavioral model alongside the static graph. The following excerpt shows
a fragment of the state-space IR for the \texttt{/users} endpoint:

\begin{Shaded}
\begin{Highlighting}[]
\FunctionTok{\{}
  \DataTypeTok{"variables"}\FunctionTok{:} \OtherTok{[}
    \FunctionTok{\{}\DataTypeTok{"name"}\FunctionTok{:} \StringTok{"request.method"}\FunctionTok{,} \DataTypeTok{"type"}\FunctionTok{:} \StringTok{"str"}\FunctionTok{,} \DataTypeTok{"domain"}\FunctionTok{:} \OtherTok{[}\StringTok{"GET"}\OtherTok{,} \StringTok{"POST"}\OtherTok{]}\FunctionTok{\}}\OtherTok{,}
    \FunctionTok{\{}\DataTypeTok{"name"}\FunctionTok{:} \StringTok{"db\_connected"}\FunctionTok{,}   \DataTypeTok{"type"}\FunctionTok{:} \StringTok{"bool"}\FunctionTok{,} \DataTypeTok{"domain"}\FunctionTok{:} \OtherTok{[}\KeywordTok{true}\OtherTok{,} \KeywordTok{false}\OtherTok{]}\FunctionTok{\}}\OtherTok{,}
    \FunctionTok{\{}\DataTypeTok{"name"}\FunctionTok{:} \StringTok{"response\_code"}\FunctionTok{,}  \DataTypeTok{"type"}\FunctionTok{:} \StringTok{"int"}\FunctionTok{,}  \DataTypeTok{"domain"}\FunctionTok{:} \OtherTok{[}\DecValTok{200}\OtherTok{,} \DecValTok{400}\OtherTok{,} \DecValTok{500}\OtherTok{]}\FunctionTok{\}}
  \OtherTok{]}\FunctionTok{,}
  \DataTypeTok{"actions"}\FunctionTok{:} \OtherTok{[}
    \FunctionTok{\{}\DataTypeTok{"name"}\FunctionTok{:} \StringTok{"validate\_input"}\FunctionTok{,} \DataTypeTok{"source"}\FunctionTok{:} \StringTok{"routes/users.py:18"}\FunctionTok{\}}\OtherTok{,}
    \FunctionTok{\{}\DataTypeTok{"name"}\FunctionTok{:} \StringTok{"query\_database"}\FunctionTok{,}  \DataTypeTok{"source"}\FunctionTok{:} \StringTok{"routes/users.py:22"}\FunctionTok{\}}\OtherTok{,}
    \FunctionTok{\{}\DataTypeTok{"name"}\FunctionTok{:} \StringTok{"format\_response"}\FunctionTok{,} \DataTypeTok{"source"}\FunctionTok{:} \StringTok{"routes/users.py:30"}\FunctionTok{\}}
  \OtherTok{]}\FunctionTok{,}
  \DataTypeTok{"transitions"}\FunctionTok{:} \OtherTok{[}
    \FunctionTok{\{}
      \DataTypeTok{"from\_state"}\FunctionTok{:} \FunctionTok{\{}\DataTypeTok{"request.method"}\FunctionTok{:} \StringTok{"GET"}\FunctionTok{,} \DataTypeTok{"db\_connected"}\FunctionTok{:} \KeywordTok{true}\FunctionTok{\},}
      \DataTypeTok{"action"}\FunctionTok{:} \StringTok{"query\_database"}\FunctionTok{,}
      \DataTypeTok{"to\_state"}\FunctionTok{:} \FunctionTok{\{}\DataTypeTok{"response\_code"}\FunctionTok{:} \DecValTok{200}\FunctionTok{\},}
      \DataTypeTok{"confidence"}\FunctionTok{:} \FloatTok{0.94}\FunctionTok{,}
      \DataTypeTok{"tier"}\FunctionTok{:} \StringTok{"STATIC\_PLUS\_RUNTIME"}
    \FunctionTok{\}}\OtherTok{,}
    \FunctionTok{\{}
      \DataTypeTok{"from\_state"}\FunctionTok{:} \FunctionTok{\{}\DataTypeTok{"request.method"}\FunctionTok{:} \StringTok{"POST"}\FunctionTok{,} \DataTypeTok{"db\_connected"}\FunctionTok{:} \KeywordTok{true}\FunctionTok{\},}
      \DataTypeTok{"action"}\FunctionTok{:} \StringTok{"validate\_input"}\FunctionTok{,}
      \DataTypeTok{"to\_state"}\FunctionTok{:} \FunctionTok{\{}\DataTypeTok{"response\_code"}\FunctionTok{:} \DecValTok{400}\FunctionTok{\},}
      \DataTypeTok{"confidence"}\FunctionTok{:} \FloatTok{0.71}\FunctionTok{,}
      \DataTypeTok{"tier"}\FunctionTok{:} \StringTok{"STATIC\_ONLY"}
    \FunctionTok{\}}\OtherTok{,}
    \FunctionTok{\{}
      \DataTypeTok{"from\_state"}\FunctionTok{:} \FunctionTok{\{}\DataTypeTok{"db\_connected"}\FunctionTok{:} \KeywordTok{false}\FunctionTok{\},}
      \DataTypeTok{"action"}\FunctionTok{:} \StringTok{"query\_database"}\FunctionTok{,}
      \DataTypeTok{"to\_state"}\FunctionTok{:} \FunctionTok{\{}\DataTypeTok{"response\_code"}\FunctionTok{:} \DecValTok{500}\FunctionTok{\},}
      \DataTypeTok{"confidence"}\FunctionTok{:} \FloatTok{0.58}\FunctionTok{,}
      \DataTypeTok{"tier"}\FunctionTok{:} \StringTok{"RUNTIME\_ONLY"}
    \FunctionTok{\}}
  \OtherTok{]}
\FunctionTok{\}}
\end{Highlighting}
\end{Shaded}

Confidence tiers follow the thresholds defined in
\texttt{determine\_confidence\_tier}: \textbf{STATIC\_PLUS\_RUNTIME}
(\(c \geq 0.65\), with both static and dynamic evidence) indicates
corroboration from AST analysis and execution traces;
\textbf{STATIC\_ONLY} (\(c \geq 0.5\), static evidence only) reflects
assertions grounded in source structure alone; \textbf{RUNTIME\_ONLY}
(\(c \geq 0.4\), dynamic evidence only) flags inferences from runtime
data without static corroboration. A fourth tier,
\textbf{HUMAN\_REVIEWED} (\(c \geq 0.9\), with human review evidence),
is available for manually curated mappings.

\subsubsection{Failure modes and graceful
degradation}\label{failure-modes-and-graceful-degradation}

COGANT is designed so that missing or partial inputs degrade the output
bundle rather than halt it, and the manuscript records these degradation
paths explicitly so that downstream consumers can interpret a partial
run without guessing what was skipped \citep{peng2011reproducible}. Five
failure modes are worth naming because each has a visible signature in
the emitted artifacts.

\textbf{No dynamic traces available.} When neither coverage data nor
execution traces are supplied, \texttt{enrich\_graph()} is a no-op: no
\texttt{coverage\_hits}, \texttt{branch\_coverage},
\texttt{call\_count}, \texttt{avg\_duration\_ms}, or
\texttt{is\_hot\_path} metadata is attached, and the graph's
\texttt{evidence\_sources} list never acquires
\texttt{dynamic\_coverage} or \texttt{dynamic\_trace} markers. The
state-space compiler still runs end-to-end using the semantic mappings
and static graph alone, but every resulting transition is confined to
the \texttt{STATIC\_ONLY} tier (or lower), and no transition can be
promoted to \texttt{STATIC\_PLUS\_RUNTIME}. Downstream code that
consumes the bundle can identify a purely-static run at a glance by
checking the evidence source markers on the graph metadata rather than
scanning individual transitions.

\textbf{Incomplete coverage data.} When coverage is supplied but covers
only a subset of the project, \texttt{\_enrich\_with\_coverage}
annotates the matched files only. The enrichment loop walks each
coverage span, normalizes the reported file path, looks up candidate
nodes whose path matches, and annotates those whose
\texttt{source\_range} overlaps the covered line; unmatched files are
silently skipped with an informational log line, and nodes in unmatched
files retain no coverage metadata at all. The result is a partially
enriched graph in which some regions are eligible for tier promotion and
others are not. Because the enrichment summary returns the number of
annotated nodes, users can check whether the observed annotation rate
matches their expectations for the provided coverage input.

\textbf{Translation rules that do not match.} When the active rule set
fails to produce any mapping for a node -- whether because no rule
fires, or because all firing rules are pruned during conflict resolution
-- that node remains outside \texttt{self.mappings} and simply does not
appear in any \texttt{graph\_fragment\_node\_ids}. The translation
engine's \texttt{get\_coverage\_report()} reports this directly: the
returned dictionary contains the total node count, the covered count, a
two-decimal \texttt{coverage\_percent}, and the sorted list of
\texttt{uncovered\_node\_ids}. Rather than treat uncovered nodes as an
error, the engine surfaces them as an explicit gap list so that authors
can either extend the rule set, hand the gap list to the
\texttt{ReviewAPI} for curation, or accept the partial mapping for
downstream Generalized Notation Notation (GNN) export.

\textbf{Fixpoint non-convergence.} The translation engine bounds its
fixpoint loop at \texttt{max\_iterations\ =\ 10} by default. If a rule
set is pathological enough to keep emitting new mappings past that bound
-- for example because two rules can produce mutually triggering
mappings -- the engine emits a
\texttt{"Max\ iterations\ reached\ without\ convergence"} warning, stops
the loop, and proceeds to conflict resolution with whatever mappings it
has accumulated. The iteration cap therefore guarantees termination even
for misconfigured rule sets, and the warning in the log gives rule
authors a clear signal that the cap was hit. Because each iteration also
writes an \texttt{iteration\_complete} entry into the engine's internal
match log, the exact per-pass mapping counts are available for diagnosis
without rerunning the pipeline.

\textbf{Pipeline error tolerance.} The default pipeline configuration
(\texttt{cogant/py/cogant/config/defaults.py}) marks the
\texttt{dynamic}, \texttt{translate}, \texttt{state\_space}, and
\texttt{process} stages with \texttt{skip\_on\_error=True}, while
structural stages such as \texttt{ingest}, \texttt{static},
\texttt{graph}, \texttt{export}, and \texttt{validate} keep the stricter
default \texttt{skip\_on\_error=False}. When an optional stage raises,
the runner records the error in the bundle, emits a warning, and
continues to the next stage rather than aborting the run. Downstream
bundle accessors (\texttt{state\_space\_model}, \texttt{process\_model})
simply return \texttt{None} for stages that did not produce output, so a
downstream consumer can detect a skipped stage by checking its accessor
rather than by parsing log text. The net effect is that a partial bundle
-- for example, a program graph and semantic mappings without a
state-space model -- remains a first-class artifact with a clear
provenance trail, rather than an opaque failure.

\subsection{Rust layer}\label{rust-layer}

Native crates (\texttt{cogant-core}, \texttt{cogant-graph},
\texttt{cogant-translate}, \texttt{cogant-statespace},
\texttt{cogant-gnn}, \texttt{cogant-ffi}, and related packages)
implement typed graph operations and export formatting. When PyO3
bindings are active, Python delegates heavy graph work through
\texttt{cogant-ffi}. Where bindings are not yet wired for a code path,
Python fallbacks apply; see SPEC \textbf{Implementation status} for the
current boundary.

\newpage

\section{Conclusion}\label{conclusion}

COGANT frames codebase analysis as a pipeline from ingestion through a
program graph IR to \textbf{Generalized Notation Notation (GNN)}
exports, with explicit confidence and provenance so that learning
systems can treat analysis noise as data rather than as hidden failure
modes. The Python layer provides session and pipeline APIs, a bundle
abstraction, CLI, review tooling, and HTML reporting; the Rust layer
concentrates graph mechanics and export formatting under crate
boundaries described in \texttt{../cogant/docs/ARCHITECTURE.md}.

\subsection{Shipped Capabilities}\label{shipped-capabilities}

Several capabilities ship in the current v0.1.x line and together define
the behaviour that downstream users can rely on:

\begin{enumerate}
\def\labelenumi{\arabic{enumi}.}
\tightlist
\item
  \textbf{Fixpoint translation engine.} The shipped
  \texttt{TranslationEngine} re-applies every registered rule on each
  pass, terminates as soon as a pass produces zero new mappings, and
  bounds pathological rule sets with a configurable iteration cap. Each
  iteration boundary is recorded in the internal match log, so
  convergence can be audited after the fact.
\item
  \textbf{Rule priority and conflict resolution.} Overlapping mappings
  are reconciled by confidence score: when two mappings cover
  overlapping \texttt{graph\_fragment\_node\_ids}, the engine retains
  the higher-confidence mapping and records a
  \texttt{conflict\_resolved} event naming the winner, the loser, and
  the overlap. Rule priority is therefore expressed through the
  confidence model rather than a separate ordering table, and
  \texttt{translate\_with\_confidence()} re-scores and re-resolves so
  that priority shifts from the \texttt{ConfidenceModel} are honoured.
\item
  \textbf{Dynamic enrichment from coverage and traces.} When
  \texttt{.coverage} SQLite or Cobertura XML inputs and Chrome DevTools
  traces are supplied, \texttt{enrich\_graph()} annotates nodes with
  \texttt{coverage\_hits}, \texttt{branch\_coverage},
  \texttt{call\_count}, \texttt{avg\_duration\_ms}, and
  \texttt{is\_hot\_path} metadata, and appends
  \texttt{dynamic\_coverage} and \texttt{dynamic\_trace} markers to the
  program graph's evidence sources.
\item
  \textbf{Confidence tier promotion.} Mappings whose evidence set
  acquires dynamic markers become eligible for promotion from
  \texttt{STATIC\_ONLY} to \texttt{STATIC\_PLUS\_RUNTIME}, and hot-path
  and branch-coverage metadata raise the underlying score through the
  diversity bonus in Equation \ref{eq:confidence-core}. A purely-static
  run remains a first-class bundle; it is simply marked as such on the
  graph metadata.
\item
  \textbf{Ten-stage pipeline architecture.} The runner executes an
  ordered DAG (ingest, static, normalize, graph, translate, state space,
  process, export, validate, and report), with
  \texttt{skip\_on\_error=True} on optional stages so that a missing
  dynamic input or a pathological rule set produces a partial but
  well-formed Generalized Notation Notation bundle rather than an opaque
  failure.
\end{enumerate}

\textbf{Limitations} follow the honest scope in
\texttt{../cogant/docs/SPEC.md}: multi-language support beyond Python is
largely roadmap; translation rules and state-space extraction are
\textbf{partial} and repository-dependent; native acceleration is
staged. Users should validate exports on their own corpora before
trusting downstream model metrics.

\textbf{Intended users} include researchers building datasets from
open-source repositories, teams prototyping Active Inference models or
graph neural network training pipelines over program-graph data who need
a single export contract, and engineers extending the system via
\texttt{../cogant/docs/PLUGIN\_API.md}.

\textbf{Validation} in the software-engineering sense is split: the
repository's verification report enumerates implemented modules and
entry points; scientific validation of model quality remains the
responsibility of downstream training and evaluation code.

\subsection{Roadmap and Future
Extensions}\label{roadmap-and-future-extensions}

Several concrete directions extend the current system:

\begin{enumerate}
\def\labelenumi{\arabic{enumi}.}
\item
  \textbf{Multi-language parsers.} The v0.1.x front end targets Python;
  adding parsers for JavaScript/TypeScript, Java, and Go would cover the
  majority of open-source ML-relevant repositories. Each parser
  implements the plugin interface documented in
  \texttt{../cogant/docs/PLUGIN\_API.md}, so language additions do not
  require changes to the core IR or export pipeline.
\item
  \textbf{Rust acceleration of critical paths.} The native crate layer
  (\texttt{cogant-graph}, \texttt{cogant-translate}) already defines
  typed graph operations. Wiring PyO3 bindings for the hot paths ---
  deduplication, rule matching, and Generalized Notation Notation
  section/tensor packing in \texttt{cogant-gnn} --- would reduce
  end-to-end latency for large repositories from minutes to seconds,
  based on preliminary profiling of the Python fallback implementations.
\item
  \textbf{Incremental re-analysis.} Currently the pipeline processes a
  full repository snapshot on each invocation. An incremental mode that
  accepts a Git diff and updates only affected subgraphs would enable
  integration into CI/CD workflows where per-commit turnaround matters.
  The stable-identifier scheme already supports cross-run node matching;
  the missing piece is a dependency tracker that invalidates downstream
  IRs when upstream nodes change.
\item
  \textbf{LLM-assisted rule discovery.} Translation rules are currently
  hand-authored. A semi-automated workflow could present a large
  language model with unannotated graph fragments and ask it to propose
  candidate rules, which a human reviewer then accepts, edits, or
  rejects via the existing \texttt{ReviewAPI}. This combines the
  pattern-recognition strength of LLMs with the auditability of
  declarative rules.
\item
  \textbf{Cross-repository graph linking.} When multiple repositories
  share interfaces (for example, a library and its consumers), linking
  their program graphs at call boundaries produces a richer training
  signal for tasks such as API misuse detection and cross-project code
  search. The graph homomorphism property defined in Section 2 provides
  the formal basis for identifying shared interface nodes across
  independently analyzed repositories.
\end{enumerate}

\newpage

\section{Experimental setup}\label{experimental-setup}

\subsection{Environment}\label{environment}

\textbf{Requirements}: Python 3.11 or newer (enforced in
\texttt{pyproject.toml}), plus an optional Rust toolchain
(\texttt{cargo}, stable 1.70+) when building native acceleration crates
under \texttt{../cogant/rust/}.

From the COGANT package root \href{../cogant/}{\texttt{../cogant/}}
(where \texttt{pyproject.toml} and \texttt{py/cogant/} live), install
with \texttt{uv\ sync\ -\/-all-extras}, or
\texttt{pip\ install\ -e\ ".{[}dev,viz{]}"} /
\texttt{pip\ install\ -e\ ".{[}all{]}"} as in
\href{../cogant/GETTING_STARTED.md}{GETTING\_STARTED.md}. Run those
commands from that directory when working inside this monorepo layout.
Python sources live under
\href{../cogant/py/cogant/}{\texttt{../cogant/py/cogant/}}; see the
package \href{../cogant/README.md}{README.md}.

\subsection{Running the API}\label{running-the-api}

Minimal \textbf{Session} run:

\begin{Shaded}
\begin{Highlighting}[]
\ImportTok{from}\NormalTok{ cogant }\ImportTok{import}\NormalTok{ Session}

\NormalTok{session }\OperatorTok{=}\NormalTok{ Session.from\_target(}\StringTok{"./path/to/repo"}\NormalTok{)}
\NormalTok{session.extract\_static()}
\NormalTok{session.build\_graph()}
\NormalTok{session.translate\_to\_gnn()}
\NormalTok{session.export\_all(}\StringTok{"output/"}\NormalTok{)}
\end{Highlighting}
\end{Shaded}

Minimal \textbf{Pipeline} run:

\begin{Shaded}
\begin{Highlighting}[]
\ImportTok{from}\NormalTok{ cogant }\ImportTok{import}\NormalTok{ PipelineRunner}
\ImportTok{from}\NormalTok{ cogant.api.pipeline }\ImportTok{import}\NormalTok{ PipelineConfig}

\NormalTok{runner }\OperatorTok{=}\NormalTok{ PipelineRunner()}
\NormalTok{config }\OperatorTok{=}\NormalTok{ PipelineConfig(output\_dir}\OperatorTok{=}\StringTok{"output/"}\NormalTok{)}
\NormalTok{bundle }\OperatorTok{=}\NormalTok{ runner.run(}\StringTok{"./path/to/repo"}\NormalTok{, config)}
\end{Highlighting}
\end{Shaded}

Adjust \texttt{PipelineConfig.stages}, \texttt{skip\_stages}, and
\texttt{plugins} to match the languages and tooling available on the
machine.

\subsection{Configuration files}\label{configuration-files}

YAML configuration can drive pipeline behavior (paths, stages, plugin
options). The architecture and SPEC documents describe the configuration
surface; keep project-specific secrets out of version control.

A minimal pipeline configuration looks like this:

\begin{Shaded}
\begin{Highlighting}[]
\CommentTok{\# cogant.config.yaml}
\FunctionTok{stages}\KeywordTok{:}
\AttributeTok{  }\KeywordTok{{-}}\AttributeTok{ ingest}
\AttributeTok{  }\KeywordTok{{-}}\AttributeTok{ static}
\AttributeTok{  }\KeywordTok{{-}}\AttributeTok{ normalize}
\AttributeTok{  }\KeywordTok{{-}}\AttributeTok{ graph}
\AttributeTok{  }\KeywordTok{{-}}\AttributeTok{ dynamic}\CommentTok{       \# optional; skipped if no coverage/trace inputs}
\AttributeTok{  }\KeywordTok{{-}}\AttributeTok{ translate}
\AttributeTok{  }\KeywordTok{{-}}\AttributeTok{ statespace}
\AttributeTok{  }\KeywordTok{{-}}\AttributeTok{ process}
\AttributeTok{  }\KeywordTok{{-}}\AttributeTok{ export}
\AttributeTok{  }\KeywordTok{{-}}\AttributeTok{ validate}

\FunctionTok{skip\_stages}\KeywordTok{:}\AttributeTok{ }\KeywordTok{[]}\CommentTok{   \# e.g. ["dynamic"] to force static{-}only runs}

\FunctionTok{plugins}\KeywordTok{:}
\AttributeTok{  }\FunctionTok{dynamic}\KeywordTok{:}
\AttributeTok{    }\FunctionTok{coverage\_path}\KeywordTok{:}\AttributeTok{ }\StringTok{"./coverage.xml"}
\AttributeTok{    }\FunctionTok{trace\_path}\KeywordTok{:}\AttributeTok{ }\StringTok{"./chrome\_trace.json"}
\AttributeTok{    }\FunctionTok{hot\_path\_percentile}\KeywordTok{:}\AttributeTok{ }\DecValTok{10}
\AttributeTok{  }\FunctionTok{translate}\KeywordTok{:}
\AttributeTok{    }\FunctionTok{max\_iterations}\KeywordTok{:}\AttributeTok{ }\DecValTok{10}
\AttributeTok{    }\FunctionTok{static\_only\_threshold}\KeywordTok{:}\AttributeTok{ }\FloatTok{0.5}
\AttributeTok{    }\FunctionTok{static\_plus\_runtime\_threshold}\KeywordTok{:}\AttributeTok{ }\FloatTok{0.65}

\FunctionTok{output\_dir}\KeywordTok{:}\AttributeTok{ }\StringTok{"./cogant\_output/"}
\FunctionTok{verbose}\KeywordTok{:}\AttributeTok{ }\CharTok{true}
\FunctionTok{dry\_run}\KeywordTok{:}\AttributeTok{ }\CharTok{false}
\end{Highlighting}
\end{Shaded}

Each stage key corresponds to a handler in
\texttt{cogant.api.pipeline.PipelineRunner.stage\_handlers}; plugin
sub-dictionaries are passed through to the stage at invocation time. The
threshold keys under \texttt{plugins.translate} mirror the constants
defined in \texttt{cogant.translate.confidence.ConfidenceModel}, keeping
the YAML surface and the Python defaults in lockstep.

\subsection{CLI}\label{cli}

Use the \texttt{cogant} CLI for scripted batch runs---see
\texttt{../cogant/docs/CLI\_GUIDE.md} for the command list that matches
the installed version.

\subsection{Export targets}\label{export-targets}

The primary export targets are the \textbf{Generalized Notation Notation
(GNN)} canonical Markdown (\texttt{model.gnn.md}) and the equivalent
companion JSON files described in
\texttt{../cogant/docs/GNN\_EXPORT.md}. Optional interop targets
(GraphML, Parquet) support analysis in Gephi/yEd and DuckDB, and
optional tensor views for PyTorch Geometric, DGL, or HDF5 can be
selected when downstream graph neural network training pipelines need to
consume the program graph as a relational tensor. Ensure the Python
environment includes optional dependencies for these tensor exports when
those code paths are used.

\subsection{Python AST parser
capabilities}\label{python-ast-parser-capabilities}

The v0.1.x front end relies on
\texttt{cogant.static.parser.PythonASTParser}, which processes Python
source through the standard-library \texttt{ast} module at the CPython
version available in the runtime (3.9+ recommended). The parser extracts
the following construct categories:

\begin{itemize}
\tightlist
\item
  \textbf{Module-level entities}: module docstrings,
  \texttt{\_\_all\_\_} exports, top-level assignments.
\item
  \textbf{Functions and methods}: \texttt{def} and \texttt{async\ def},
  including signatures with positional, keyword, variadic
  (\texttt{*args}, \texttt{**kwargs}), and positional-only parameters.
  Default values are recorded as constant expressions where statically
  evaluable.
\item
  \textbf{Classes}: class definitions, base classes, metaclasses, and
  the \texttt{\_\_init\_\_} / \texttt{\_\_new\_\_} boundary.
\item
  \textbf{Decorators}: \texttt{@staticmethod}, \texttt{@classmethod},
  \texttt{@property}, \texttt{@dataclass}, and arbitrary user-defined
  decorators. Decorator arguments are captured as attribute metadata.
\item
  \textbf{Type annotations}: PEP 484 / 526 / 604 annotations on function
  parameters, return types, and variable assignments. Generic subscripts
  (\texttt{List{[}int{]}}, \texttt{Dict{[}str,\ Any{]}}) are preserved
  as type strings.
\item
  \textbf{Comprehensions and generators}: list, set, dict comprehensions
  and generator expressions are represented as anonymous FUNCTION nodes
  with DATA\_FLOW edges to their enclosing scope.
\item
  \textbf{Control flow}: \texttt{if}/\texttt{elif}/\texttt{else},
  \texttt{for}/\texttt{while}/\texttt{else},
  \texttt{try}/\texttt{except}/\texttt{finally},
  \texttt{with}/\texttt{async\ with}, and \texttt{match}/\texttt{case}
  (Python 3.10+) are mapped to CONTROLFLOW\_NODE entities.
\item
  \textbf{Imports}: \texttt{import} and \texttt{from\ ...\ import}
  statements produce MODULE\_IMPORT roles with edges to the resolved
  module when discoverable on the file system.
\item
  \textbf{Constants}: module-level and class-level assignments to
  \texttt{Final} or ALL\_CAPS names are classified as CONSTANT nodes.
\end{itemize}

Constructs that require runtime evaluation (for example \texttt{exec},
\texttt{importlib.import\_module}, or dynamic \texttt{\_\_getattr\_\_})
are recorded as EXTERNAL nodes with HEURISTIC provenance and
correspondingly lower confidence.

\subsection{Progressive IR stages}\label{progressive-ir-stages}

Processing advances through six intermediate representations, each
adding semantic detail atop its predecessor. Table 3 summarizes what
each stage contributes.

\textbf{Table 3. Progressive IR stages and their contributions.}

{\def\LTcaptype{none} % do not increment counter
\begin{longtable}[]{@{}
  >{\raggedright\arraybackslash}p{(\linewidth - 6\tabcolsep) * \real{0.0986}}
  >{\raggedright\arraybackslash}p{(\linewidth - 6\tabcolsep) * \real{0.1268}}
  >{\raggedright\arraybackslash}p{(\linewidth - 6\tabcolsep) * \real{0.2113}}
  >{\raggedright\arraybackslash}p{(\linewidth - 6\tabcolsep) * \real{0.5634}}@{}}
\toprule\noalign{}
\begin{minipage}[b]{\linewidth}\raggedright
Stage
\end{minipage} & \begin{minipage}[b]{\linewidth}\raggedright
IR name
\end{minipage} & \begin{minipage}[b]{\linewidth}\raggedright
Key additions
\end{minipage} & \begin{minipage}[b]{\linewidth}\raggedright
Typical output size (10K-function repo)
\end{minipage} \\
\midrule\noalign{}
\endhead
\bottomrule\noalign{}
\endlastfoot
1 & Repo IR & Raw entities and relationships per file; deduplication;
merged type info & \textasciitilde15 MB JSON \\
2 & Program Graph IR & Consolidated directed graph \(G=(V,E)\); stable
identifiers; confidence and provenance on every node and edge &
\textasciitilde20 MB JSON \\
3 & Semantic Mapping IR & Translation rules applied; semantic roles
assigned; confidence adjusted by rule engine & \textasciitilde22 MB JSON
(graph + mapping log) \\
4 & State Space IR & Variables, actions, transitions, observations;
dynamic traces integrated where available & \textasciitilde5 MB JSON
(behavioral model) \\
5 & Process Model IR & Higher-level control patterns (request--response,
producer--consumer, state machines) & \textasciitilde2 MB JSON \\
6 & Validation IR & Coverage metrics, confidence distribution, schema
compliance, consistency checks, reproducibility hashes &
\textasciitilde1 MB JSON (report) \\
\end{longtable}
}

Stages 4 and 5 are \textbf{partial} for many repositories: the
state-space compiler requires either execution traces or sufficient
static structure (for example annotated state machines) to produce
meaningful output. The pipeline tolerates missing stages gracefully; the
Validation IR records which stages completed and which were skipped.

\subsection{Performance
characteristics}\label{performance-characteristics}

The architecture targets the following benchmarks on a 4-core machine,
as specified in \texttt{../cogant/docs/ARCHITECTURE.md}:

{\def\LTcaptype{none} % do not increment counter
\begin{longtable}[]{@{}lll@{}}
\toprule\noalign{}
Repository size & Target wall-clock time & Memory budget \\
\midrule\noalign{}
\endhead
\bottomrule\noalign{}
\endlastfoot
10K functions & \textless{} 30 s & \textless{} 500 MB \\
100K functions & \textless{} 5 min & \textless{} 2 GB \\
1M functions & \textless{} 1 hr & \textless{} 2 GB (streaming) \\
\end{longtable}
}

These targets assume the Python orchestration layer with Rust
acceleration on critical paths (graph construction, rule matching, and
Generalized Notation Notation section/tensor packing in
\texttt{cogant-gnn}). In the current v0.1.x release, where Rust bindings
are staged rather than fully wired, Python fallback implementations
handle graph operations; users should expect roughly 2--3\(\times\)
slower throughput on repositories above 50K functions until native
acceleration is enabled.

Parallelism is exploited at multiple stages: file parsing is
embarrassingly parallel, translation rules are applied concurrently over
independent subgraphs, and export formatting runs in parallel per output
format. Incremental caching skips unchanged stages when the input hash
matches a prior run, which is particularly effective during iterative
development against a single repository.

\subsection{What to record}\label{what-to-record}

For reproducible experiments, record: COGANT version or commit hash,
interpreter version, list of stages executed, configuration file
contents (redacted), input repository commit hash, and random seeds for
any learned components \textbf{outside} COGANT that consume the exports.

\newpage

\section{Reproducibility}\label{reproducibility}

Reproducible computational research requires version-pinned tools, fixed
inputs, and documented outputs \citep{peng2011reproducible}. COGANT
contributes the \textbf{graph-generation} slice of that story: identical
source trees and identical pipeline configuration should yield identical
exported bundles modulo declared nondeterminism (for example optional
neural embeddings if enabled).

\subsection{Version pinning}\label{version-pinning}

Pin:

\begin{itemize}
\tightlist
\item
  The \textbf{COGANT} version (\texttt{pyproject.toml} / package
  \texttt{\_\_version\_\_}) and Git commit when installing from source.
\item
  The \textbf{Python} minor version.
\item
  \textbf{Optional} Rust toolchain commit when building native
  extensions from \texttt{../cogant/rust/}.
\end{itemize}

\subsection{Artifact layout}\label{artifact-layout}

Pipeline output typically includes JSON for intermediate IRs,
\textbf{Generalized Notation Notation (GNN)} bundles (canonical
\texttt{model.gnn.md} plus companion JSON and interop artifacts),
validation reports, and optional HTML sites. Paths under the chosen
\texttt{output\_dir} should be treated as disposable but
\textbf{checksumable}: store hashes alongside published datasets derived
from COGANT exports when releasing research artifacts.

\subsection{Determinism}\label{determinism}

Parsing and graph construction aim for deterministic ordering on a fixed
filesystem snapshot. Features that pull in external models (for example
optional name or documentation embeddings consumed by the Generalized
Notation Notation exporter) introduce variability unless models and
seeds are fixed; the Generalized Notation Notation export document
(\texttt{../cogant/docs/GNN\_EXPORT.md}) calls out embedding dimensions
and optional behavior.

\subsection{Relation to the template
repository}\label{relation-to-the-template-repository}

While this project remains outside
\href{../../../projects/}{\texttt{../../../projects/}}, it is
\textbf{not} executed by the root \texttt{./run.sh} discovery layer.
After promotion to
\href{../../../projects/cogant/}{\texttt{../../../projects/cogant/}}
with \texttt{src/}, \texttt{tests/}, and \texttt{pyproject.toml} per
template rules, the standard manuscript validation and PDF stages apply;
until then, validate Markdown from the template repository root,
e.g.~\texttt{uv\ run\ python\ -m\ infrastructure.validation.cli\ markdown\ ./projects\_in\_progress/cogant/manuscript/}.

\subsection{Validation gates}\label{validation-gates}

The pipeline enforces quality through three complementary validation
checkers, each targeting a different failure mode:

\textbf{IntegrityChecker.} Verifies structural soundness of the program
graph: all edge endpoints reference existing nodes, no unintended
duplicate nodes exist, orphaned nodes (zero in-degree and zero
out-degree) are flagged, and self-loops are reported unless explicitly
allowed by configuration. The checker also ensures that confidence
scores fall within \([0, 1]\) and that provenance records are non-empty
for every node and edge. A graph that fails integrity checks receives a
FAIL validation status; downstream export is blocked.

\textbf{SchemaChecker.} Validates each IR artifact against its JSON
schema (versioned alongside the COGANT package). Schema violations --
such as missing required fields, incorrect types, or unknown enum values
in \texttt{NodeKind} or \texttt{SemanticRole} -- are classified as ERROR
or FATAL depending on severity. Schema versions are recorded in the
validation report so that consumers can verify compatibility with their
import code.

\textbf{ProvenanceChecker.} Audits the provenance chain: every assertion
in the semantic mapping must trace back to at least one evidence source
(SourceCode, TypeSystem, ControlFlow, Heuristic, or External). The
checker flags mappings whose provenance is empty or whose confidence
score is inconsistent with the declared evidence tier -- for example, a
STATIC\_PLUS\_RUNTIME tier with no runtime trace evidence. These flags
appear as warnings rather than errors, since partial provenance is
expected for heuristic rules.

Together, these gates ensure that exported bundles meet a minimum
quality bar before reaching downstream models. The thresholds are
configurable (see \texttt{../cogant/docs/VALIDATION.md}); defaults
require \(\geq 95\%\) coverage, mean confidence \(\geq 0.85\), and zero
schema violations.

\subsection{Pipeline checkpoint and resume
behavior}\label{pipeline-checkpoint-and-resume-behavior}

The pipeline writes a checkpoint file after each successfully completed
stage. The checkpoint records the stage name, completion timestamp,
input hash, and output hash. On a subsequent invocation with
\texttt{-\/-resume} (or \texttt{PipelineConfig.resume\ =\ True}), the
runner reads the checkpoint file, verifies that the input hash for each
completed stage still matches the current source tree, and skips stages
whose outputs are already valid. If a source file has changed, the
runner invalidates the affected stage and all downstream stages, then
re-executes from the earliest invalidated point.

Checkpoints are stored as JSON under the configured
\texttt{output\_dir}:

\begin{verbatim}
output/
  .cogant_checkpoint.json
\end{verbatim}

This mechanism is particularly useful during iterative development:
after an initial full run, changing a single source file triggers only
re-parsing, graph construction, and downstream stages for the affected
subgraph, rather than re-analyzing the entire repository.

\subsection{Output manifest structure}\label{output-manifest-structure}

Each pipeline run produces a manifest file (\texttt{manifest.json})
alongside the exported artifacts. The manifest provides a
machine-readable inventory of everything the run produced:

\begin{Shaded}
\begin{Highlighting}[]
\FunctionTok{\{}
  \DataTypeTok{"cogant\_version"}\FunctionTok{:} \StringTok{"0.1.0"}\FunctionTok{,}
  \DataTypeTok{"run\_id"}\FunctionTok{:} \StringTok{"run\_20260408\_143012"}\FunctionTok{,}
  \DataTypeTok{"timestamp"}\FunctionTok{:} \StringTok{"2026{-}04{-}08T14:30:12Z"}\FunctionTok{,}
  \DataTypeTok{"input\_hash"}\FunctionTok{:} \StringTok{"sha256:a1b2c3..."}\FunctionTok{,}
  \DataTypeTok{"config\_hash"}\FunctionTok{:} \StringTok{"sha256:d4e5f6..."}\FunctionTok{,}
  \DataTypeTok{"stages\_completed"}\FunctionTok{:} \OtherTok{[}\StringTok{"discovery"}\OtherTok{,} \StringTok{"parsing"}\OtherTok{,} \StringTok{"repo\_ir"}\OtherTok{,} \StringTok{"graph"}\OtherTok{,}
                        \StringTok{"translate"}\OtherTok{,} \StringTok{"validate"}\OtherTok{,} \StringTok{"export"}\OtherTok{]}\FunctionTok{,}
  \DataTypeTok{"stages\_skipped"}\FunctionTok{:} \OtherTok{[}\StringTok{"statespace"}\OtherTok{,} \StringTok{"process"}\OtherTok{]}\FunctionTok{,}
  \DataTypeTok{"artifacts"}\FunctionTok{:} \OtherTok{[}
    \FunctionTok{\{}\DataTypeTok{"path"}\FunctionTok{:} \StringTok{"model.gnn.json"}\FunctionTok{,} \DataTypeTok{"format"}\FunctionTok{:} \StringTok{"json"}\FunctionTok{,} \DataTypeTok{"size\_bytes"}\FunctionTok{:} \DecValTok{2048576}\FunctionTok{,}
     \DataTypeTok{"hash"}\FunctionTok{:} \StringTok{"sha256:..."}\FunctionTok{\}}\OtherTok{,}
    \FunctionTok{\{}\DataTypeTok{"path"}\FunctionTok{:} \StringTok{"model.pyg.pt"}\FunctionTok{,}   \DataTypeTok{"format"}\FunctionTok{:} \StringTok{"pytorch\_geometric"}\FunctionTok{,} \DataTypeTok{"size\_bytes"}\FunctionTok{:} \DecValTok{1024000}\FunctionTok{,}
     \DataTypeTok{"hash"}\FunctionTok{:} \StringTok{"sha256:..."}\FunctionTok{\}}\OtherTok{,}
    \FunctionTok{\{}\DataTypeTok{"path"}\FunctionTok{:} \StringTok{"validation\_report.json"}\FunctionTok{,} \DataTypeTok{"format"}\FunctionTok{:} \StringTok{"json"}\FunctionTok{,} \DataTypeTok{"size\_bytes"}\FunctionTok{:} \DecValTok{8192}\FunctionTok{,}
     \DataTypeTok{"hash"}\FunctionTok{:} \StringTok{"sha256:..."}\FunctionTok{\}}
  \OtherTok{]}\FunctionTok{,}
  \DataTypeTok{"statistics"}\FunctionTok{:} \FunctionTok{\{}
    \DataTypeTok{"nodes"}\FunctionTok{:} \DecValTok{320}\FunctionTok{,}
    \DataTypeTok{"edges"}\FunctionTok{:} \DecValTok{1250}\FunctionTok{,}
    \DataTypeTok{"mean\_confidence"}\FunctionTok{:} \FloatTok{0.92}\FunctionTok{,}
    \DataTypeTok{"validation\_status"}\FunctionTok{:} \StringTok{"PASS"}
  \FunctionTok{\}}
\FunctionTok{\}}
\end{Highlighting}
\end{Shaded}

Storing per-artifact SHA-256 hashes enables consumers to verify that
exported bundles have not been modified after generation. When
publishing datasets derived from COGANT exports, including the manifest
alongside the data satisfies the reproducibility requirements outlined
by \citep{peng2011reproducible}: any researcher with the same source
tree, COGANT version, and configuration can regenerate the artifacts and
compare hashes.

\subsection{Data ethics and licensing}\label{data-ethics-and-licensing}

Exported graphs can contain identifiers and comments from source code.
Redistribution of derived graphs must respect the licenses of input
repositories and organizational data policies.

\newpage

\section{Scope and related work}\label{scope-and-related-work}

\subsection{Positioning}\label{positioning}

\textbf{Machine learning for big code} surveys cover naturalness,
representation learning, and task taxonomy \citep{allamanis2018survey}.
COGANT contributes \textbf{infrastructure}: a documented IR, pipeline,
and export contract rather than a new benchmark or model architecture.

\textbf{Graph neural networks} unify message-passing frameworks for
relational data \citep{wu2020comprehensive, scarselli2009graph}.
Although COGANT's primary output is the Active Inference Institute's
Generalized Notation Notation, its program graph can be materialised as
optional tensor views compatible with these frameworks (PyG, DGL) so
that graph-neural-network model papers can cite a specific tensor
layout.

\textbf{Code property graphs} and related security-oriented graph
constructions show how to merge syntactic and semantic edges for
analysis \citep{yamaguchi2014modeling}. COGANT's program graph is
cognate but emphasizes \textbf{ML export}, confidence scoring, and
multi-stage IR refinement.

\subsection{Related tool categories}\label{related-tool-categories}

\begin{itemize}
\tightlist
\item
  \textbf{Compiler and LLVM IRs} --- richer semantics, heavier
  toolchain; COGANT favors lightweight extraction for dataset building.
\item
  \textbf{SCA and linters} --- rule enforcement focus; COGANT focuses on
  graph generation for downstream learning.
\item
  \textbf{Neural code models} (code2vec, CodeBERT, etc.) --- often
  consume tokens or AST paths; COGANT supplies \textbf{explicit program
  graphs} that graph-neural-network-centric research can ingest
  directly.
\end{itemize}

\subsection{Program analysis for ML}\label{program-analysis-for-ml}

Several systems address the intersection of program analysis and machine
learning that COGANT operates in:

\textbf{code2seq and code2vec} \citep{alon2019code2vec} represent
programs as sets of AST paths between terminal nodes. These path-based
representations are effective for method naming and code summarization
but discard the full graph topology that graph-neural-network-based
models exploit. COGANT preserves the complete program graph ---
including control-flow and data-flow edges --- and can export it in
tensor form, enabling downstream architectures that reason over richer
relational structure.

\textbf{Gated Graph Neural Networks (GGNN)} \citep{li2016gated}
introduced gated recurrent propagation for graph-structured program
representations, demonstrating strong results on variable misuse
detection and other code tasks. COGANT's export contract (PyG
\texttt{Data} objects with typed \texttt{edge\_index} and
\texttt{edge\_attr}) is directly compatible with GGNN-style models; the
kind and role indices provide the discrete node and edge types that
typed message-passing layers require.

\textbf{CodeQL} (Semmle/GitHub) provides a declarative query language
over relational representations of code. While CodeQL excels at security
analysis with hand-written queries, its outputs are query results rather
than tensor-ready graph bundles. COGANT occupies the complementary
niche: it produces the graph data that learned models consume, and can
ingest CodeQL query results as an additional evidence source feeding the
confidence model.

\textbf{CodeBERT} \citep{feng2020codebert} and related pre-trained
models operate at the token level, learning representations from natural
language and code jointly. These models are complementary to COGANT's
graph-centric approach: CodeBERT embeddings could serve as optional node
features in COGANT's Generalized Notation Notation (GNN) export (the
export schema already reserves dimensions for text embeddings as
documented in \texttt{../cogant/docs/GNN\_EXPORT.md}).

\subsubsection{Feature matrix: COGANT vs.~related
tools}\label{feature-matrix-cogant-vs.-related-tools}

The following matrix contrasts COGANT's capabilities with the related
tools discussed in this section. Entries marked ``✓'' indicate
first-class support; ``partial'' indicates limited or indirect support;
``---'' indicates the feature is out of scope for that tool.

\textbf{Table 4. Feature comparison of program-to-model toolchains.}

{\def\LTcaptype{none} % do not increment counter
\begin{longtable}[]{@{}lccccc@{}}
\toprule\noalign{}
Feature & COGANT & code2vec & GGNN & CodeQL & CodeBERT \\
\midrule\noalign{}
\endhead
\bottomrule\noalign{}
\endlastfoot
Full program graph (AST + CFG + DFG) & ✓ & --- & input-only & ✓ & --- \\
Typed node/edge taxonomy & ✓ & --- & partial & ✓ & --- \\
Confidence scoring per assertion & ✓ & --- & --- & --- & --- \\
Provenance tracking & ✓ & --- & --- & partial & --- \\
State-space extraction & ✓ & --- & --- & --- & --- \\
Temporal regime classification & ✓ & --- & --- & --- & --- \\
Dynamic enrichment (coverage, traces) & ✓ & --- & --- & partial & --- \\
Generalized Notation Notation output & ✓ & --- & --- & --- & --- \\
Tensor export (PyG, DGL, HDF5) & ✓ & partial & input-only & --- & --- \\
Pluggable translation rules & ✓ & --- & --- & ✓ & --- \\
Human review loop & ✓ & --- & --- & partial & --- \\
Multi-language front-ends & roadmap & ✓ & --- & ✓ & ✓ \\
\end{longtable}
}

COGANT is distinct from the other toolchains in three ways: first, it
explicitly models uncertainty through confidence tiers tied to evidence
provenance; second, it produces a structured Active Inference notation
as its primary output rather than an opaque tensor; and third, it
composes static and dynamic evidence in a single pipeline rather than
specializing to one.

\subsection{Active inference and program
behavior}\label{active-inference-and-program-behavior}

The state-space IR in COGANT's pipeline (states, actions, transitions,
observations) shares structural parallels with \textbf{active inference}
formulations \citep{friston2010free}, where an agent maintains beliefs
about hidden states and selects actions to minimize prediction error. In
the program analysis context, the ``agent'' is the analysis pipeline
itself: it observes code artifacts, maintains beliefs about program
behavior (the state-space model), and refines those beliefs as new
evidence (dynamic traces, coverage data) arrives.

This connection is analogical: the \texttt{ConfidenceModel} in
\texttt{../cogant/py/cogant/translate/confidence.py} aggregates evidence
and penalties in a way that suggests belief revision, but it is not a
Bayesian posterior. Future work could formalize a tighter link by
casting rule application as variational inference, where a fixpoint
would represent an approximate posterior over program semantics.

\subsection{Boundaries}\label{boundaries}

COGANT does not subsume formal verification, interactive theorem
proving, or full interprocedural pointer analysis unless implemented as
explicit future stages. The SPEC marks Rust acceleration and additional
parsers as staged; the manuscript should be read together with that
table for up-to-date scope.

\subsection{Forward compatibility}\label{forward-compatibility}

Promoting COGANT into
\href{../../../projects/}{\texttt{../../../projects/}} integrates
manuscript PDF rendering with the template's validation gates.
Cross-references in this folder use paths \textbf{relative to these
Markdown files} (for example
\href{../cogant/docs/}{\texttt{../cogant/docs/}}) so links stay stable
when the tree moves.

\newpage

\section{Manuscript Syntax Reference
(COGANT)}\label{manuscript-syntax-reference-cogant}

Formatting conventions for Markdown in this folder. For the full
rendering contract, see the template exemplar
\href{../../../projects/code_project/manuscript/SYNTAX.md}{\texttt{projects/code\_project/manuscript/SYNTAX.md}}.

\subsection{Citations}\label{citations}

Use Pandoc cite syntax; keys must exist in \texttt{references.bib}.

\begin{Shaded}
\begin{Highlighting}[]
\CommentTok{[}\OtherTok{@allamanis2018survey}\CommentTok{]}

\CommentTok{[}\OtherTok{@allamanis2018survey; @wu2020comprehensive}\CommentTok{]}
\end{Highlighting}
\end{Shaded}

\subsection{Equations}\label{equations}

Use LaTeX \texttt{equation} with \texttt{\textbackslash{}label} /
\texttt{\textbackslash{}ref} as documented in the code\_project SYNTAX,
or pandoc-crossref attributes where the pipeline enables them.

\subsection{Figures}\label{figures}

If you add figures, place assets where the future project
\texttt{output/} layout can resolve them (typically \texttt{../figures/}
after promotion to \texttt{projects/cogant/}). Use explicit relative
paths from the rendering contract described in
\texttt{infrastructure/rendering/AGENTS.md}.

\subsection{Section files}\label{section-files}

Numeric prefixes \texttt{00\_}--\texttt{09\_} are combined in
stem-sorted order by
\texttt{infrastructure/rendering/manuscript\_discovery.py}. Files named
\texttt{SYNTAX.md} sort in the \textbf{other} bucket (after main
sections, before \texttt{99\_*} if present).

\subsection{Cross-references to the
package}\label{cross-references-to-the-package}

Prefer \textbf{relative} paths from this folder to the package tree,
e.g.~\href{../cogant/docs/ARCHITECTURE.md}{\texttt{../cogant/docs/ARCHITECTURE.md}},
so links work in the editor and in Git without hard-coding the monorepo
path.

\subsection{See also}\label{see-also}

\begin{itemize}
\tightlist
\item
  \href{AGENTS.md}{\texttt{AGENTS.md}} --- editor protocol for this
  folder
\item
  \href{README.md}{\texttt{README.md}} --- orientation and render notes
\end{itemize}



\clearpage

\bibliography{references}
\end{document}
